\chapter*{Abstract}	
Network Intrusion has been on the rise ever since the invention of internet and IoT devices that specifically accumulate huge amounts of data. 
Network Intrusion Detection System (NIDS) has been a major topic considering the number of threats involved in today’s cyber security scenario. 
With the ever-rising threats in communication networks, development of NIDS plays a boosting factor to prevent the networks from attacks of all kinds. 
The purpose of the development of network intrusion detection system is to protect the communication network, preserve and safeguard vital information 
from intruders dynamically by analyzing the complex and sophisticated traffic data. The motives for the intrusion vary on individual and institutional 
levels, financial gains, unauthorized access of intellectual properties, confidential information, political gains and others. NIDS monitors and analyzes 
all the incoming web traffics, detects them as ‘Attack’ or ‘Benign’ thus safeguarding the network from future attacks of same kinds. Several research has 
been carried out to detect network intrusion for specific fields (Health, Finance, Entertainment, Military, Government) incorporating advanced techniques 
with the utilization of specific data. Previously published research papers showed that the deployment of robust machine learning algorithm achieved greater 
performance in intrusion detection when trained and tested on the same dataset. Research portals such as JSRT Journal, MDPI, ResearchGate, ScienceDirect, 
Springer Nature, IEEE contain research papers related to cross dataset generalization (mostly between CICIDS2017 and CSE-CIC-IDS2018). Most of the research 
papers drew their conclusions based on ‘Accuracy’ of the models for intrusion detection. There was no uniformity in these papers. Some described the usage of 
SMOTE while training CICIDS2017 dataset in supervised learning models, variation on the size of dataset and percentage of data used for training the models; 
thus, leaving a gap for cross-dataset generalization to the unseen network traffic. This research paper tries to showcase the capabilities of Supervised Learning 
Models (Random Forest, XGBoost), Unsupervised Learning Models (PCA, K-Means, Autoencoders) and Deep Learning Models in cross-dataset generalization. All the models 
will be trained on CICIDS2017 dataset and tested on two different test datasets namely CSE-CIC-IDS2018 and UNSW-NB15; which are publicly available. No SMOTE will be 
applied to oversample the CICIDS2017 dataset. Unlike previous papers, this research will prioritize the evaluation metrics significant for cross-dataset generalization. 
Finally, the performance of the models will allow us to assess the strength and practicality of a unique approach in real-world cybersecurity scenario.